\documentclass[a4paper,12pt,twoside]{ctexart}
\usepackage[dvipsnames]{xcolor}
\usepackage{geometry}
\usepackage{tcolorbox}
\usepackage{tikz}
\usepackage{amssymb}
\usepackage{graphicx}

% 定义公函标准炭黑
\definecolor{inkblack}{rgb}{0.08,0.08,0.08}

\geometry{a4paper, left=3.18cm, right=3.18cm, top=2.54cm, bottom=2.54cm}

% 配置页脚
\usepackage{fancyhdr}
\pagestyle{fancy}
\fancyhf{}
\renewcommand{\headrulewidth}{0pt}
\fancyfoot[RO]{\makebox[0pt][l]{\hspace*{28.2pt}\raisebox{2.81pt}{\rmfamily\fontsize{10.5bp}{10.5bp}\selectfont\thepage}}}
\fancyfoot[LE]{\makebox[0pt][r]{\hspace*{33.45pt}\raisebox{2.81pt}{\rmfamily\fontsize{10.5bp}{10.5bp}\selectfont\thepage}}}

% 字体
\newCJKfontfamily\fzdbs[Path=./]{FZDBS.ttf}
\newCJKfontfamily\fzssk[Path=./]{FZSSK.ttf}

\AtBeginDocument{\color{inkblack}}

% 标题
\newcommand{\lawtitle}[2]{
    \vspace*{0.507cm} 
    \begin{center}
        {\fzdbs\fontsize{22bp}{22bp}\selectfont #1}\par
        \vspace{10.37bp}
        {\fzdbs\fontsize{18bp}{18bp}\selectfont #2}\par
    \end{center}
    \vspace{18.4bp}
}

% 表框
\newtcolorbox{notesection}[1][]{
    width=467.72bp,
    sharp corners,
    colback=white,
    colframe=inkblack,
    boxrule=0.5pt,
    fontupper=\color{inkblack}\fontsize{10.5bp}{10.5bp}\selectfont,
    enlarge left by=-19.28bp,
    before=\vspace*{-13.39bp}\noindent,
    after=\par,
    boxsep=0pt,
    left=4bp,
    top=4.69bp,
    bottom=4.69bp,
    right=0pt,
    #1
}

\newcommand{\lawpar}{\par\vspace*{-4.86bp}}

% 副标题
\newcommand{\subtitlesection}[1]{%
    \begin{tcolorbox}[
        width=467.72bp,
        sharp corners,
        colback=white,
        colframe=inkblack,
        boxrule=0.5pt,
        boxsep=0pt,
        left=0pt, right=0pt,
        top=4bp, bottom=4bp,
        halign=center,
        fontupper=\fzdbs\fontsize{15bp}{15bp}\selectfont\color{inkblack},
        enlarge left by=-19.28bp,
        before={\nointerlineskip\vspace{-0.5pt}\noindent},
        after=\vspace{0pt}
    ]
    #1
    \end{tcolorbox}%
}

% 左右分栏
\newcommand{\pairsection}[2]{%
    \begin{tcolorbox}[
        width=467.72bp,
        sharp corners,
        colback=white,
        colframe=inkblack,
        boxrule=0.5pt,
        enlarge left by=-19.28bp,
        before={\nointerlineskip\vspace{-0.5pt}\noindent},
        after=\par,
        boxsep=0pt,
        left=0pt, right=0pt,
        top=0pt, bottom=0pt
    ]
    \begin{minipage}[c]{112.9bp}%
        \centering\songti\fontsize{10.5bp}{10.5bp}\selectfont #1%
    \end{minipage}%
    \hspace{0pt}%
    \vrule width 0.5pt%
    \hspace{4bp}%
    \begin{minipage}[c]{346.32bp}%
        \vspace{4bp}%
        \songti\mdseries\fontsize{10.5bp}{17bp}\selectfont%
        \baselineskip=17bp%
        \setlength{\parindent}{0pt}%
        \color{inkblack}%
        #2%
        \vspace{4bp}%
    \end{minipage}%
    \end{tcolorbox}%
}

\begin{document}

% ===== 第 1 页 =====

\lawtitle{民事起诉状}{（买卖合同纠纷）}

% 说明框
\begin{notesection}
\hfuzz=50pt
\color{inkblack}\heiti\fontsize{10.5bp}{16bp}\selectfont 说明：\par
\vspace*{-0.2bp}
\songti\mdseries\fontsize{10.5bp}{16bp}\selectfont
\setlength{\parindent}{20.7bp}
\sloppy

\setlength{\parskip}{0pt}
\vspace*{-5.1bp}为了方便您更好地参加诉讼，保护您的合法权利，请填写本表。
\lawpar
1.起诉时需向人民法院提交证明您身份的材料，如身份证复印件、营业执照复印件等。
\lawpar
2.本表所列内容是您参加诉讼以及人民法院查明案件事实所需，请务必如实填写。
\lawpar
3.本表有些内容可能与您的案件无关，您认为与案件无关的项目可以填"无"或不填；对于本\\[-4.86bp]
表中勾选项可以在对应项打"√"；您认为另有重要内容需要列明的，可以另附页填写。
\lawpar
4.本表 word 电子版填写时，相关栏目可复制粘贴或扩容，但不得改变要素内容、格式设置。例\\[-4.86bp]
如，多原告、多被告或多委托诉讼代理人等情况，可根据实际情况复制粘贴；需填写文字较多时，\\[-4.86bp]
可根据实际对栏目进行扩容等。
\lawpar
$\bigstar$特别提示$\bigstar$\\[-4.86bp]
\hspace*{20.7bp}诉讼参加人应遵守诚信原则如实认真填写表格。\\[-4.86bp]
\hspace*{20.7bp}如果诉讼参加人违反有关规定，虚假诉讼、恶意诉讼、滥用诉权，人民法院将视违法情形依法\\[-4.86bp]
追究责任。
\end{notesection}

% 当事人信息
\subtitlesection{当事人信息}

% 原告（自然人）
\pairsection{原告\\（自然人）}{
姓名：\\
性别：男$\square$\hspace{1em}女$\square$\\
出生日期：\hspace{2em}年\hspace{2em}月\hspace{2em}日\hspace{4em}民族：\\
工作单位：\hspace{8em}职务：\hspace{6em}联系电话：\\
住所地（户籍所在地）：\\
经常居住地：\\
证件类型：\\
证件号码：
}

% 原告（法人、非法人组织）
\pairsection{原告\\（法人、非法人组织）}{
名称：\\
住所地（主要办事机构所在地）：\\
注册地 / 登记地：\\
法定代表人 / 负责人：\hspace{4em}职务：\hspace{4em}联系电话：\\
统一社会信用代码：\\
类型：有限责任公司$\square$\hspace{1em}股份有限公司$\square$\hspace{1em}上市公司$\square$\\
\hspace*{36bp}其他企业法人$\square$\hspace{1em}事业单位$\square$\hspace{1em}社会团体$\square$\hspace{1em}基金会$\square$\\
\hspace*{36bp}社会服务机构$\square$\hspace{1em}机关法人$\square$\hspace{1em}农村集体经济组织法人$\square$\\
\hspace*{36bp}城镇农村的合作经济组织法人$\square$\hspace{1em}基层群众性自治组织法人$\square$\\
\hspace*{36bp}个人独资企业$\square$\hspace{1em}合伙企业$\square$\hspace{1em}不具有法人资格的专业服务机构$\square$\\
所有制性质：国有$\square$（控股$\square$\hspace{1em}参股$\square$）\hspace{1em}民营$\square$\hspace{1em}其他\rule{60bp}{0.5pt}
}

\newpage

% ===== 第 2 页 =====

% 委托诉讼代理人
\pairsection{委托诉讼代理人}{
有$\square$\\
\hspace*{21bp}姓名：\\
\hspace*{21bp}\makebox[0.33\linewidth][l]{单位：}\makebox[0.33\linewidth][l]{职务：}\makebox[0.33\linewidth][l]{联系电话：}\\
\hspace*{21bp}代理权限：一般授权$\square$\hspace{1em}特别授权$\square$\hspace{1em}\rule{80bp}{0.5pt}\\
无$\square$
}

% 被告（自然人）
\pairsection{被告\\（自然人）}{
姓名：\\
性别：男$\square$\hspace{1em}女$\square$\\
出生日期：\hspace{2em}年\hspace{2em}月\hspace{2em}日\hspace{4em}民族：\\
工作单位：\hspace{8em}职务：\hspace{6em}联系电话：\\
住所地（户籍所在地）：\\
经常居住地：\\
证件类型：\\
证件号码：
}

% 被告（法人、非法人组织）
\pairsection{被告\\（法人、非法人组织）}{
名称：\\
住所地（主要办事机构所在地）：\\
注册地 / 登记地：\\
法定代表人 / 负责人：\hspace{4em}职务：\hspace{4em}联系电话：\\
统一社会信用代码：\\
类型：有限责任公司$\square$\hspace{1em}股份有限公司$\square$\hspace{1em}上市公司$\square$\\
\hspace*{36bp}其他企业法人$\square$\hspace{1em}事业单位$\square$\hspace{1em}社会团体$\square$\hspace{1em}基金会$\square$\\
\hspace*{36bp}社会服务机构$\square$\hspace{1em}机关法人$\square$\hspace{1em}农村集体经济组织法人$\square$\\
\hspace*{36bp}城镇农村的合作经济组织法人$\square$\hspace{1em}基层群众性自治组织法人$\square$\\
\hspace*{36bp}个人独资企业$\square$\hspace{1em}合伙企业$\square$\hspace{1em}不具有法人资格的专业服务机构$\square$\\
所有制性质：国有$\square$（控股$\square$\hspace{1em}参股$\square$）\hspace{1em}民营$\square$\hspace{1em}其他\rule{60bp}{0.5pt}
}

% 第三人（自然人）
\pairsection{第三人\\（自然人）}{
姓名：\\
性别：男$\square$\hspace{1em}女$\square$\\
出生日期：\hspace{2em}年\hspace{2em}月\hspace{2em}日\hspace{4em}民族：\\
工作单位：\hspace{8em}职务：\hspace{6em}联系电话：\\
住所地（户籍所在地）：\\
经常居住地：\\
证件类型：\\
证件号码：
}

% 第三人（法人、非法人组织）— 简短版
\pairsection{第三人\\（法人、非法人组织）}{
名称：\\
住所地（主要办事机构所在地）：\\
注册地 / 登记地：\\
法定代表人 / 负责人：\hspace{4em}职务：\hspace{4em}联系电话：\\
统一社会信用代码：
}

\newpage

% ===== 第 3 页 =====

% 第三人（法人、非法人组织）— 类型部分续
\pairsection{第三人\\（法人、非法人组织）}{
类型：有限责任公司$\square$\hspace{1em}股份有限公司$\square$\hspace{1em}上市公司$\square$\\
\hspace*{36bp}其他企业法人$\square$\hspace{1em}事业单位$\square$\hspace{1em}社会团体$\square$\hspace{1em}基金会$\square$\\
\hspace*{36bp}社会服务机构$\square$\hspace{1em}机关法人$\square$\hspace{1em}农村集体经济组织法人$\square$\\
\hspace*{36bp}城镇农村的合作经济组织法人$\square$\hspace{1em}基层群众性自治组织法人$\square$\\
\hspace*{36bp}个人独资企业$\square$\hspace{1em}合伙企业$\square$\hspace{1em}不具有法人资格的专业服务机构$\square$\\
所有制性质：国有$\square$（控股$\square$\hspace{1em}参股$\square$）\hspace{1em}民营$\square$\hspace{1em}其他\rule{60bp}{0.5pt}
}

% 诉讼请求副标题（含提示文字）
\begin{tcolorbox}[
    width=467.72bp,
    sharp corners,
    colback=white,
    colframe=inkblack,
    boxrule=0.5pt,
    boxsep=0pt,
    left=0pt, right=0pt,
    top=4bp, bottom=4bp,
    halign=center,
    enlarge left by=-19.28bp,
    before={\nointerlineskip\vspace{-0.5pt}\noindent},
    after=\vspace{0pt}
]
{\fzdbs\fontsize{15bp}{15bp}\selectfont\color{inkblack}诉讼请求}\par\nointerlineskip
\vspace{2bp}
{\heiti\fontsize{10.5bp}{14bp}\selectfont\color{inkblack}（原告为卖方时，填写第 1 项、第 2 项；原告为买方时，填写第 3 项、第 4 项；\\第 5 项至第 10 项为共同项）}
\end{tcolorbox}

% 诉讼请求总述框
\begin{tcolorbox}[
    width=467.72bp,
    sharp corners,
    colback=white,
    colframe=inkblack,
    boxrule=0.5pt,
    enlarge left by=-19.28bp,
    before={\nointerlineskip\vspace{-0.5pt}\noindent},
    after=\par,
    boxsep=0pt,
    left=4bp, right=0pt,
    top=4bp, bottom=4bp
]
\songti\mdseries\fontsize{10.5bp}{17bp}\selectfont%
\baselineskip=17bp%
\parskip=0pt%
\lineskiplimit=-\maxdimen%
\setlength{\parindent}{0pt}%
\color{inkblack}%
（可完整表述诉讼请求；为方便、准确梳理要点，相关内容请在下方要素式表格中填写）\\[51bp]
\mbox{}
\end{tcolorbox}

% 1. 给付价款（元）
\pairsection{1. 给付价款（元）}{
\rule{60bp}{0.5pt}\hspace{1em}元（人民币，下同；如外币需特别注明）
}

% 2. 迟延给付价款的利息（违约金）
\pairsection{2. 迟延给付价款的利息\\（违约金）}{
截至\hspace{2em}年\hspace{2em}月\hspace{2em}日止，迟延给付价款的利息\hspace{4em}元、违约金\\
\hspace{2em}元，自\hspace{4em}之后的逾期利息、违约金，以\hspace{4em}元为基数\\
按照\hspace{4em}标准计算；\\
计算方式：\\
是否请求支付至实际清偿之日止：是$\square$\hspace{1em}否$\square$
}

% 3. 赔偿因卖方违约所受的损失
\pairsection{3. 赔偿因卖方违约所受\\的损失}{
支付赔偿金\hspace{4em}元\\
违约类型：迟延履行$\square$\hspace{1em}不履行$\square$\hspace{1em}其他$\square$\\
具体情形：\\
损失计算依据：
}

% 4. 是否对标的物的瑕疵承担责任
\pairsection{4. 是否对标的物的瑕疵\\承担责任}{
是$\square$\hspace{1em}修理$\square$\hspace{1em}重作$\square$\hspace{1em}更换$\square$\hspace{1em}退货$\square$\hspace{1em}减少价款或者报酬$\square$\\
\hspace*{21bp}其他$\square$：\\
否$\square$
}

% 5. 要求继续履行或者解除合同
\pairsection{5. 要求继续履行或者解\\除合同}{
继续履行$\square$\hspace{1em}日内履行完毕付款$\square$\hspace{1em}供货$\square$\hspace{1em}义务\\
判令解除合同$\square$\\
确认买卖合同已于\hspace{2em}年\hspace{2em}月\hspace{2em}日解除$\square$
}

% 6. 是否主张担保权利
\pairsection{6. 是否主张担保权利}{
是$\square$\hspace{4em}内容：\\
否$\square$
}

% 7. 是否主张实现债权的费用
\pairsection{7. 是否主张实现债权的\\费用}{
是$\square$\hspace{4em}费用明细：\\
否$\square$
}

% 8. 是否主张诉讼费用
\pairsection{8. 是否主张诉讼费用}{
是$\square$\\
否$\square$
}

% 9. 其他请求
\pairsection{9. 其他请求}{
\mbox{}
}

% 10. 标的总额
\pairsection{10. 标的总额}{
\mbox{}
}

\newpage

% ===== 第 4 页 =====

% 约定管辖和诉前保全
\subtitlesection{约定管辖和诉前保全}

% 1. 有无仲裁、法院管辖约定
\pairsection{1. 有无仲裁、法院管辖\\约定}{
有$\square$\hspace{4em}合同条款及内容：\\
无$\square$
}

% 2. 是否已经诉前保全
\pairsection{2. 是否已经诉前保全}{
是$\square$\hspace{1em}保全法院：\hspace{4em}保全时间：\\
\hspace*{21bp}保全案号：\\
否$\square$\\
（如申请诉讼保全，请另行提交诉讼保全申请及相关材料）
}

% 事实与理由
\subtitlesection{事实与理由}

% 事实与理由总述框
\begin{tcolorbox}[
    width=467.72bp,
    sharp corners,
    colback=white,
    colframe=inkblack,
    boxrule=0.5pt,
    enlarge left by=-19.28bp,
    before={\nointerlineskip\vspace{-0.5pt}\noindent},
    after=\par,
    boxsep=0pt,
    left=4bp, right=0pt,
    top=4bp, bottom=4bp
]
\songti\mdseries\fontsize{10.5bp}{17bp}\selectfont%
\baselineskip=17bp%
\parskip=0pt%
\lineskiplimit=-\maxdimen%
\setlength{\parindent}{0pt}%
\color{inkblack}%
（可完整表述纠纷涉及的事实与理由；为方便、准确梳理要点，相关内容请在下方要素式表格中填写）\\[51bp]
\mbox{}
\end{tcolorbox}

% 1. 合同的签订情况
\pairsection{1. 合同的签订情况\\（名称、编号、签订时\\间、地点等；如无书\\面合同，请注明"无\\书面合同"）}{
\mbox{}
}

% 2. 合同主体
\pairsection{2. 合同主体}{
出卖人（卖方）：\\
买受人（买方）：
}

% 3. 买卖标的物情况
\pairsection{3. 买卖标的物情况\\（标的物名称、规格、\\质量、数量等）}{
\mbox{}
}

% 4. 合同约定的价格及支付方式
\pairsection{4. 合同约定的价格及支\\付方式}{
单价\hspace{4em}元；总价\hspace{4em}元；\\
以现金$\square$\hspace{1em}转账$\square$\hspace{1em}票据$\square$（写明票据类型）\hspace{1em}其他$\square$\hspace{1em}方式\\
一次性$\square$\hspace{1em}分期$\square$\hspace{1em}支付\\
分期方式：
}

% 5. 合同约定的交货时间、地点、方式、风险承担、安装、调试、验收
\pairsection{5. 合同约定的交货时\\间、地点、方式、风\\险承担、安装、调试、\\验收}{
\mbox{}
}

% 6. 合同约定的质量标准及检验方式、质量异议期限
\pairsection{6. 合同约定的质量标准\\及检验方式、质量异\\议期限}{
\mbox{}
}

% 7. 合同约定的违约金（定金）
\pairsection{7. 合同约定的违约金\\（定金）}{
违约金$\square$\hspace{4em}元（合同条款：第\hspace{2em}条）\\
定金$\square$\hspace{4em}元（合同条款：第\hspace{2em}条）\\
迟延履行违约金$\square$\hspace{2em}\%/ 日（合同条款：第\hspace{2em}条）
}

\newpage

% ===== 第 5 页 =====

% 8. 价款支付及标的物交付情况
\pairsection{8. 价款支付及标的物交\\付情况}{
按期支付价款\hspace{4em}元，逾期付款\hspace{4em}元，逾期未付款\hspace{4em}元\\
按期交付标的物\hspace{4em}件，逾期交付\hspace{4em}件，逾期未交付\hspace{4em}件
}

% 9. 是否存在迟延履行
\pairsection{9. 是否存在迟延履行}{
是$\square$\hspace{1em}迟延时间：\hspace{6em}逾期付款$\square$\hspace{1em}逾期交货$\square$\\
否$\square$
}

% 10. 是否催促过履行
\pairsection{10. 是否催促过履行}{
是$\square$\hspace{1em}催促情况：\hspace{2em}年\hspace{2em}月\hspace{2em}日通过\hspace{4em}方式进行了催促\\
否$\square$
}

% 11. 买卖合同标的物有无质量争议
\pairsection{11. 买卖合同标的物有\\无质量争议}{
有$\square$\hspace{1em}具体情况：\\
无$\square$
}

% 12. 标的物质量规格或履行方式是否存在不符合约定的情况
\pairsection{12. 标的物质量规格或\\履行方式是否存在不\\符合约定的情况}{
是$\square$\hspace{1em}具体情况：\\
否$\square$
}

% 13. 是否曾就标的物质量问题进行协商
\pairsection{13. 是否曾就标的物质\\量问题进行协商}{
是$\square$\hspace{1em}具体情况：\\
否$\square$
}

% 14. 是否通知解除合同
\pairsection{14. 是否通知解除合同}{
是$\square$\hspace{1em}具体情况：\\
否$\square$
}

% 15. 被告应当支付的利息、违约金、赔偿金
\pairsection{15. 被告应当支付的利\\息、违约金、赔偿金}{
利息$\square$\hspace{4em}元\\
违约金$\square$\hspace{4em}元\\
赔偿金$\square$\hspace{4em}元\\
共计\hspace{4em}元\hspace{4em}计算方式：
}

% 16. 是否签订物的担保（抵押、质押）合同
\pairsection{16. 是否签订物的担保\\（抵押、质押）合同}{
是$\square$\hspace{1em}签订时间：\\
否$\square$
}

% 17. 担保人、担保物
\pairsection{17. 担保人、担保物}{
担保人：\\
担保物：
}

% 18. 是否最高额担保（抵押、质押）
\pairsection{18. 是否最高额担保\\（抵押、质押）}{
是$\square$\hspace{1em}担保债权的确定时间：\\
\hspace*{21bp}担保额度：\\
否$\square$
}

% 19. 是否办理抵押、质押登记
\pairsection{19. 是否办理抵押、质\\押登记}{
是$\square$\hspace{1em}正式登记$\square$\\
\hspace*{21bp}预告登记$\square$\\
否$\square$
}

% 20. 是否签订保证合同
\pairsection{20. 是否签订保证合同}{
是$\square$\hspace{1em}签订时间：\hspace{4em}保证人：\hspace{4em}主要内容：\\
否$\square$
}

% 21. 保证方式
\pairsection{21. 保证方式}{
一般保证$\square$\\
连带责任保证$\square$
}

% 22. 其他担保方式
\pairsection{22. 其他担保方式}{
是$\square$\hspace{1em}形式：\\
否$\square$
}

% 23. 请求承担责任的依据
\pairsection{23. 请求承担责任的依据}{
合同约定：\\
法律规定：
}

\newpage

% ===== 第 6 页 =====

% 24. 其他需要说明的内容
\pairsection{24. 其他需要说明的内\\容（可另附页）}{
\mbox{}
}

% 25. 证据清单
\pairsection{25. 证据清单（可另附\\页）}{
\mbox{}
}

% 对纠纷解决方式的意愿
\subtitlesection{对纠纷解决方式的意愿}

% 是否了解调解作为非诉讼纠纷解决方式
\pairsection{是否了解调解作为非诉\\讼纠纷解决方式，能\\及时、高效、低成本、\\不伤和气地解决纠纷}{
了解$\square$\hspace{1em}不了解$\square$
}

% 是否了解先行调解解决纠纷的好处
\pairsection{是否了解先行调解解\\决纠纷的好处}{
1. 立案后选择先行调解的，可以很快启动调解程序。如不同意调解，法院\\
将依程序开庭审理案件，但可能需要经过较长一段时间的排期等待，且审\\
理、执行周期相对较长。\\
了解$\square$\hspace{1em}不了解$\square$\\
2. 选择先行调解，调解成功且自动履行的免交诉讼费用，申请司法确认的\\
不交纳诉讼费用，要求出具调解书的减半交纳诉讼费用。\\
了解$\square$\hspace{1em}不了解$\square$\\
3. 首次调解不成功，但仍有继续调解意愿的，可以选择更换调解组织和调\\
解员再进行调解。调解无法达成一致意见的，法院将依程序排期开庭。\\
了解$\square$\hspace{1em}不了解$\square$\\
4. 依照法律规定，调解具有保密性要求，调解过程不公开，调解协议未经\\
当事人同意不得公开。\\
了解$\square$\hspace{1em}不了解$\square$\\
5. 调解达成的协议具有法律效力，可以依照法律规定申请司法确认，具有\\
强制执行效力。\\
了解$\square$\hspace{1em}不了解$\square$
}

% 是否考虑先行调解
\pairsection{是否考虑先行调解}{
是$\square$\\
否$\square$\\
暂不确定，想要了解更多内容$\square$
}

% 具状人签字与日期
\vspace{10bp}
\noindent\hspace*{248bp}{\fzdbs\fontsize{15bp}{22bp}\selectfont 具状人（签字、盖章）：}\par
\noindent\hspace*{248bp}{\fzdbs\fontsize{15bp}{22bp}\selectfont 日期：}

\end{document}
